\section{Complete classification of the residual square class thirteen}
\label{q13-section}

The following closes one actual nonunit residual class. It is independent
of the modular candidate enumeration and does not assert uniformity over
moving residual classes or exponents. The rank conclusion below is proved
by explicit descent; a preliminary software rank calculation is not an
input to the proof.

\begin{lemma}[The required rational elliptic group]\label{q13-group}
For
\[
 E:\quad Y^2=X^3-13X^2-507X
\]
one has $E(\mathbb Q)=\{O,(0,0)\}$.
\end{lemma}
\begin{proof}
Its rational two-isogenous curve is
\[
 E':\quad Y^2=X^3+26X^2+2197X.
\]
Use the explicit descent of \eqref{qg-isogeny}--\eqref{qg-covering-quartic}.
For $Y^2=X^3+AX^2+BX$, the homomorphism to rational square classes
has values $1$ at $O$, $B$ at $(0,0)$, and $X$ elsewhere.
Its kernel is the image of the dual isogeny. Every nonexceptional
image class is represented by a signed square-free divisor $d$ of $B$
with a cover
\begin{equation}\label{q13-cover}
 N^2=dU^4+AU^2V^2+(B/d)V^4,\qquad \gcd(U,V)=1.
\end{equation}
These statements include the exceptional points: the dual preimage of
$(0,0)$ exists exactly when $B$ is a square. They follow from the
explicit isogeny and substitution $X=d(U/V)^2$, as in the earlier
ordinary descent proof.

For $E$, the candidate classes are $\pm1,\pm3,\pm13,\pm39$.
The classes $1,-3$ occur at $O,(0,0)$. Multiplication by $-3$ pairs
the other classes, so it suffices to exclude $-1,13,-13$.
For $d=-1$, the cover is
\[
 N^2=-U^4-13U^2V^2+507V^4.
\]
Coprimality leaves the following exhaustive possibilities modulo $16$:
\[
\begin{array}{c|c}
 \text{parity}&\text{right side modulo }16\\\hline
 U\text{ odd},\ V\text{ even}&15,11\\
 U\text{ even},\ V\text{ odd}&11,7\\
 U,V\text{ odd}&13,5
\end{array}
\]
None is a square residue.

For $d=13$, the cover gives
$N^2=13(U^4-U^2V^2-3V^4)$, so $13\mid N$ and
$U^4-U^2V^2-3V^4\equiv0\pmod{13}$. If $13\mid V$, then also
$13\mid U$, a contradiction. Otherwise $z=(U/V)^2$ modulo $13$
satisfies
\[
 z^2-z-3=(z-7)^2\equiv0\pmod{13},
\]
but $7$ is not a square modulo $13$. The class $d=-13$ similarly
forces $z^2+z-3=(z-6)^2\equiv0$, and $6$ is not a square either.
Thus the image is exactly $\alpha(E)=\{1,-3\}$.

On $E'$, the quadratic factor
$X^2+26X+2197=(X+13)^2+2028$ is positive, so every affine
point with $X\ne0$ has positive $X$. The signed-divisor test for $2197=13^3$
leaves only classes $1,13$, both attained by $O$ and
$T'=(0,0)$. Hence $\alpha'(E')=\{1,13\}$.
Writing the isogenies as $\phi:E\to E'$ and $\psi:E'\to E$,
every $P'\in E'(\mathbb Q)$ has either $P'$ or $P'+T'$ in
$\ker\alpha'=\phi(E(\mathbb Q))$. Since $\psi(T')=O$ and
$\psi\phi=[2]$, this proves
\[
 \psi(E'(\mathbb Q))=2E(\mathbb Q),\qquad
 \ker\alpha=2E(\mathbb Q),\qquad
 [E(\mathbb Q):2E(\mathbb Q)]=2.
\]
There is exactly one nonzero rational point of order two: the other
quadratic factor has discriminant $2197$, a nonsquare. Mordell--Weil
finite generation therefore identifies this index as $2^{\operatorname{rank}E+1}$,
giving rank zero.

The discriminant $16\cdot507^2\cdot2197$ has support in $\{2,3,13\}$.
At the good primes $5,7$, complete counts of affine ordinates at
successive abscissas are
\[
 (1,2,0,2,0),\qquad (1,0,0,0,0,2,0).
\]
Including infinity gives orders $6$ and $4$. Good-reduction injection
on prime-to-residue-characteristic torsion shows that the rational
torsion order divides $2$: its $5$-primary part is excluded at $7$,
its $7$-primary part at $5$, and every other primary part divides both
orders. The known point $(0,0)$ gives equality. This proves the lemma.
The finite-generation and torsion-injection inputs are the standard
ones already specified in the preceding descent section.
\end{proof}

\begin{lemma}[An integral map to the thirteen twist]\label{q13-integral-map}
For arbitrary integers $a,b,s$, put
\[
 c=a-b,\quad y=a+b,\quad A=7a^2+12ab+7b^2,
 \qquad U=13(A-26s),\quad V=13Uy.
\]
Then
\begin{align}
 A^2+3c^4&=52F(a,b),&c^2+13y^2&=2A,\label{q13-integral-basic}\\
 V^2-U^3+13U^2c^2+507Uc^4
   &=8788U\bigl(F(a,b)-13s^2\bigr).\label{q13-integral-cubic}
\end{align}
If $F(a,b)=13s^2$ and $c\ne0$, then $U\ne0$, and
\[
 (X,Y)=(U/c^2,V/c^3)
\]
is a finite rational point of $E$ with $X\ne0$.
\end{lemma}
\begin{proof}
The first two identities follow by expansion. Factoring the left side
of \eqref{q13-integral-cubic} by $U$, its bracket becomes
\[
 169Uy^2-U^2+13Uc^2+507c^4
 =26AU-U^2+507c^4
 =169(A^2-676s^2+3c^4),
\]
which is $8788(F-13s^2)$ by \eqref{q13-integral-basic}.
If $U=0$, then $A=26s$, and the first identity with $F=13s^2$
forces $3c^4=0$, a contradiction. Dividing the resulting homogeneous
cubic identity by $c^6$ proves the curve equation. The only denominator
condition is $c\ne0$; there is no division by $a,b,s,A$ or $V$.
\end{proof}

\begin{theorem}[The actual thirteen-square class]\label{q13-classification}
Let $a,b$ be coprime positive integers and let $s\in\mathbb Z$.
Then
\[
 F(a,b)=13s^2
 \quad\Longleftrightarrow\quad
 a=b=1\ \text{ and }\ s\in\{-1,1\}.
\]
Consequently no such seed has $F=13Q^g$ with $Q>1$ and a positive
even exponent $g$.
\end{theorem}
\begin{proof}
If $a\ne b$, Lemma~\ref{q13-integral-map} gives a finite point of
$E$ with nonzero $X$, contradicting Lemma~\ref{q13-group}.
Thus $a=b$, and positivity and coprimality give $a=b=1$.
Substitution gives $s^2=1$. The converse is direct. For the final
assertion put $s=Q^{g/2}$; the classification gives $Q^{g/2}=1$,
contrary to $Q>1$.
\end{proof}

\begin{corollary}[The homogeneous integer classification]\label{q13-all-integers}
For arbitrary integers $a,b,s$,
\[
 F(a,b)=13s^2
 \quad\Longleftrightarrow\quad a=b\ \text{ and }\ s\in\{-a^2,a^2\}.
\]
\end{corollary}
\begin{proof}
The integral map and rational group classification impose $a=b$
without positivity or coprimality. Substitution gives $s^2=a^4$, hence
$(s-a^2)(s+a^2)=0$. The converse follows by the same substitution,
including $a=b=s=0$.
\end{proof}

\paragraph{Verification scope.}
The three local descent exclusions, both good-prime point counts and
the explicit rational map were independently reviewed. The integral
map is a separately reviewed simplification whose sole denominator is
$a-b$. Exact replay checks the polynomial identities and all admissible
residue pairs and the original rational-map identities; it does not
supply the elliptic group theorem by a rank database call. Scoped Lean
verification of the integer identities and nonvanishing verifies that
bridge, not Mordell--Weil, the two-isogeny descent, or the complete
rational-point classification. The final Lean diagonal statement has
the integral-point obstruction as an explicit hypothesis.
